\begin{problem}{}{}

    Demostrar las siguientes afirmaciones, donde $a$, $b$ son siempre números enteros.
    Justificar cada uno de los pasos en cada demostración indicando el axioma o resultado que utiliza.

    \begin{enumerate}[(a)]
        \item $a + a = a$ implica que $a = 0$.
        \item $(a + b) \cdot c = a \cdot c + b \cdot c$.
        \item $-a = (-1) \cdot a$.
    \end{enumerate}

\end{problem}
\vspace{1ex}

Vamos a ver cómo, partiendo desde los axiomas dados en clase y explicitando cada paso, podemos demostrar cada una de las proposiciones dadas:

\begin{enumerate}[(a)]
    \item La idea intuitiva radica en usar el axioma de inverso aditivo hasta lograr eliminar las ``a'' del lado izquierdo. Como sigue: \\
        \begin{tabular}{|l l r}
             & $a + a = a$ & \\
            $\implies$ & $a + a + (-a) = a + (-a)$ & (propiedad uniforme y existencia del inverso aditivo) \\
            $\iff$ & $a + (a + (-a)) = (a + (-a))$ & (propiedad asociativa de la suma) \\
            $\iff$ & $a + 0 = 0$ & (Leibniz) \\
            $\iff$ & $a = 0$ & (elemento neutro de la suma y Leibniz) \\
        \end{tabular}

        De la propiedad transitiva de la implicación, se deduce la afirmación deseada.

    \item Analizando la expresión, este es una versión del axioma de distributividad en orden inverso.
          Por esto, podemos partir de este y usar los axiomas de conmutatividad para acomodar el resultado.
          Como sigue: \\
          \begin{tabular}{|l l r}
             & $c \cdot (a + b) = c \cdot a + c \cdot b$ & (axioma de distributividad) \\
             $\iff$ & $(a + b) \cdot c = a \cdot c + b \cdot c$ & (conmutatividad y Leibniz) \\
          \end{tabular}

          De esta forma vemos que la identidad dada es siempre verdadera.

    \item Finalmente, podemos pensar la igualdad como \textit{``El inverso aditivo de $a$ es $(-1)\cdot a$''}.
          Para esto, podemos probar que la suma de $(-1) \cdot a$ con $a$ es nula. Como sigue: \\
          \begin{tabular}{|l l r}
             & $a + (-1) \cdot a$ & \\
             $=$ & $1 \cdot a + (-1) \cdot a$ & (neutro multiplicativo) \\
             $=$ & $(1 + -1) \cdot a$ & (inciso (b)) \\
             $=$ & $0 \cdot a$ & (inverso aditivo y Leibniz) \\
             $=$ & $0$ & (elemento absorvente del producto) \\
          \end{tabular}

          De este último, necesariamente $(-1) \cdot a$ es el inverso multiplicativo de $a$ y por unicidad del opuesto es igual a $-a$.
\end{enumerate}
